% ============================================================
% f2013_zhejiang_science_q12_views.tex
% 2013 浙江卷（理科）第 12 题 三视图（TikZ 矢量重绘）
% 正视图：宽 4、高 3+2=5 的矩形，斜线自左边距底 2 处连至右上顶点；
% 侧视图：宽 3、高 5 的矩形，斜线自左上顶点连至右边距底 2 处；
% 俯视图：两直角边 4、3 的直角三角形，斜边为上方对角线。
% ============================================================
\begin{tikzpicture}[
  x=1cm, y=1cm, line width=0.5pt, font=\footnotesize, >=Stealth
]

  % ---------- 正视图 ----------
  \draw (0,0) rectangle (4,5);
  \draw (0,2) -- (4,5);
  \draw[|<->|] (0,5.40) -- (4,5.40) node[midway, above] {$4$};
  \foreach \x in {0,4} \draw (\x,5.22) -- (\x,5.58);
  \node at (2,-0.62) {正视图};

  % ---------- 正视图右侧的高度标注（3 与 2） ----------
  \draw[|<->|] (4.52,5) -- (4.52,2) node[midway, right] {$3$};
  \draw[|<->|] (4.52,2) -- (4.52,0) node[midway, right] {$2$};
  \foreach \y in {0,2,5} \draw (4.34,\y) -- (4.70,\y);

  % ---------- 侧视图 ----------
  \draw (5.80,0) rectangle (8.80,5);
  \draw (5.80,5) -- (8.80,2);
  \draw[|<->|] (5.80,5.40) -- (8.80,5.40) node[midway, above] {$3$};
  \foreach \x in {5.80,8.80} \draw (\x,5.22) -- (\x,5.58);
  \node at (7.30,-0.62) {侧视图};

  % ---------- 俯视图 ----------
  \draw (0,-1.50) -- (0,-4.50) -- (4,-4.50) -- cycle;
  \draw[|<->|] (4.52,-1.50) -- (4.52,-4.50) node[midway, right] {$3$};
  \foreach \y in {-1.50,-4.50} \draw (4.34,\y) -- (4.70,\y);
  \node at (2,-5.02) {俯视图};

\end{tikzpicture}
